% This is a modified version of Springer's LNCS template suitable for anonymized MICCAI 2025 main conference submissions. 
% Original file: samplepaper.tex, a sample chapter demonstrating the LLNCS macro package for Springer Computer Science proceedings; Version 2.21 of 2022/01/12

\documentclass[runningheads]{llncs}
%
\usepackage[T1]{fontenc}
\usepackage{algorithm}
\usepackage{algorithmic}
\usepackage{amsmath,amssymb}
% T1 fonts will be used to generate the final print and online PDFs,
% so please use T1 fonts in your manuscript whenever possible.
% Other font encodings may result in incorrect characters.
%
\usepackage{graphicx,verbatim}
% Used for displaying a sample figure. If possible, figure files should
% be included in EPS format.
%
% If you use the hyperref package, please uncomment the following two lines
% to display URLs in blue roman font according to Springer's eBook style:
%\usepackage{color}
%\renewcommand\UrlFont{\color{blue}\rmfamily}
%\urlstyle{rm}
%
\begin{document}
%
\title{Privacy-Preserving Federated Learning for Imbalanced Medical Applications via Distribution-Aware Aggregation}
%\titlerunning{Abbreviated paper title}
% If the paper title is too long for the running head, you can set
% an abbreviated paper title here
%
\begin{comment}  %% Removed for anonymized MICCAI submission
\author{First Author\inst{1}\orcidID{0000-1111-2222-3333} \and
Second Author\inst{2,3}\orcidID{1111-2222-3333-4444} \and
Third Author\inst{3}\orcidID{2222--3333-4444-5555}}
%
\authorrunning{F. Author et al.}
% First names are abbreviated in the running head.
% If there are more than two authors, 'et al.' is used.
%
\institute{Princeton University, Princeton NJ 08544, USA \and
Springer Heidelberg, Tiergartenstr. 17, 69121 Heidelberg, Germany
\email{lncs@springer.com}\\
\url{http://www.springer.com/gp/computer-science/lncs} \and
ABC Institute, Rupert-Karls-University Heidelberg, Heidelberg, Germany\\
\email{\{abc,lncs\}@uni-heidelberg.de}}

\end{comment}

\author{Anonymized Authors}  %% Added for anonymized MICCAI submission
\authorrunning{Anonymized Author et al.}
\institute{Anonymized Affiliations \\
    \email{email@anonymized.com}}
  
\maketitle              % typeset the header of the contribution
%
\begin{abstract}
Federated Learning (FL) enables collaborative model training across institutions without sharing raw data. However, performance often degrades when local datasets exhibit imbalanced label distributions, a challenge inherent to medical settings where healthy or common cases dominate over rare pathologies across sites. Existing approaches to label skew often rely on model-specific aggregation or auxiliary training objectives, assuming prior knowledge of global label distributions, potentially weakening privacy guarantees.
To address this, we propose a general-purpose FL framework that dynamically adapts its aggregation strategy based on a privacy-preserving estimate of the global label distribution. This distribution is estimated using a lightweight federated analytics protocol combining secure aggregation with differential privacy. As strong privacy guarantees inevitably introduce noise, we explicitly account for this effect by modeling the uncertainty associated with the estimated global label distribution.
Federated training begins with standard FL aggregation and transitions to a distribution-aware mechanism once a reliable global label distribution estimation becomes available. We evaluate the framework on medical imaging and physiological signal classification tasks under realistic heterogeneous settings. Results demonstrate consistent improvements over classical FL aggregation methods in non-IID label-skew scenarios, with negligible resource overhead and minimal loss attributed to the privacy-preserving mechanism. Our approach offers a practical, model-agnostic solution for federated medical applications where both data imbalance and privacy are critical concerns.

\keywords{Federated Learning  \and NON-IID Data Distribution \and Privacy.}
% Authors must provide keywords and are not allowed to remove this Keyword section.

\end{abstract}
%
%
%
\section{Important Messages to MICCAI Authors}

This is a modified version of Springer's LNCS template suitable for anonymized MICCAI 2026 main conference submissions. The author section has already been anonymized for you.  You cannot remove the author section to gain extra writing space.  You must not remove the abstract and keyword sections in your submission.

To avoid an accidental anonymity breach, you are not required to include the \textbf{Acknowledgments} and the \textbf{Disclosure of Interest} sections for the initial submission phase.  If your paper is accepted, you must reinsert these sections in your final paper.  If you opted to include these sections in the initial submission phase, they will count toward the 8-page limit.

To avoid desk rejection of your paper, you are encouraged to read "Avoiding Desk Rejection" by the MICCAI submission platform manager.  A few important rules are summarized below for your reference:  
\begin{enumerate}
    \item Do not modify or remove the provided anonymized author section except for inserting your paper ID.
    \item Do not remove the keyword section to gain extra writing space. Authors must provide at least one keyword.
    \item Do not remove the abstract sections to gain extra writing space. 
    \item Inline figures (wrapping text around a figure) are not allowed.
    \item Authors are not allowed to change the default margins, font size, font type, and document style.  If you are using any additional packages, it is the author's responsibility to ensure they do not inherently alter the document's layout.
    \item It is prohibited to use commands such as \verb|\vspace| and \verb|\hspace| to reduce the pre-set white space in the document.  Shrinking the white space between your figures/tables and their captions, and between sections and subsections, is strictly prohibited.
    \item Please make sure your figures and tables do not span into the margins.
    \item If you have tables, the smallest font size allowed is 8pt.  Your tables should not be inserted as figures.
    \item Please ensure your paper is properly anonymized.  Refer to your previously published paper in the third person.  Data collection location of private dataset should be masked.
        \item  If you are including a link to your code repository or datasets or submitting supplementary materials, you must ensure your repository and supplementary materials are anonymized.
    \item Your paper must not exceed the page limit.  You can have 8 pages of the main content (text (including paper title, author section, abstract, keyword section), figures and tables, conclusion, acknowledgment, and disclaimer) plus up to 2 pages of references.  The reference section does not need to start on a new page.  
    \item *Beginning in 2025* 
    Supplementary materials are limited to multimedia content (e.g., videos) as warranted by the technical application (e.g., robotics, surgery, ...). These files should not display any proofs, analysis, additional results, or embedded slides, and should not show any identification markers either. Violation of this guideline will lead to desk rejection. PDF files may not be submitted as supplementary materials in 2026 unless authors are citing a paper that has not yet been published. In such a case, authors are required to submit an anonymized version of the cited paper. ACs will perform checks to ensure the integrity.
\end{enumerate}
\section{Related Work Papers}
\textbf{FedALA: Adaptive Local Aggregation for Personalized Federated Learning}
\\\textbf{Adaptive Personalized Federated Learning (UPenn)}
\\\textbf{Federated Learning with Label Distribution Skew via Logits Calibration}
\\Client Side Adaptive FL (altra tecnica per fare adaptive federated learning client-side)
\\\textbf{FedMcon: an adaptive aggregation method for federated learning via meta controller}
\\\textbf{Federated Learning with Label Distribution Skew via Logits Calibration (FedLC)}
\\requires local data to train a controller that manages aggregation weights assignment
\\giustuficazioen del label e quantity skew : \textbf{Federated Learning on Non-IID Data Silos: An Experimental Study}

\section{Introduction}


\section{Related Work}
\textbf{Federated Learning with non-IID data}   In Real-World FL applications, one of the main challenges is keeping the model's performance and convergence speed consistent when data is unevenly distributed among clients. This scenario is generally referred to as a non-IID data distribution, and it includes different types of data heterogeneity. Non-IID distribution scenarios are mainly connected to two phenomena. Statistical heterogeneity, where feature distributions vary across clients, and label heterogeneity, where class proportions or label availability differ among clients. Both of these settings significantly affect training stability and global model generalization.

\textbf{Model Related Methods} To mitigate label non-IID effects in federated learning, most approaches go beyond the FedAvg paradigm by explicitly using model related feature to weight the updates. Methods such as SCaFFOLD, FedProx, and FedNova introduce model-related control variables or regularization terms designed to align client updates and reduce client drift, with the goal of stabilizing the optimization process and preventing divergence of the global model. However, a key limitation of these approaches lies in their limited adaptability across different model architectures, as their effectiveness often depends on architectural assumptions or specific optimization behaviors.

More recent approaches, including FedDual and FedALA, address data heterogeneity by redistributing the optimization load between the server and the clients. By granting greater flexibility to the clients through partial, adaptive, or decoupled updates, these methods can improve training stability under highly heterogeneous settings. Nevertheless, this increased client-side freedom introduces a critical trade-off: locally relevant information may not be fully captured by the global aggregation process, leading to the loss of fine-grained details that are not effectively incorporated into the global model and subsequently redistributed. As a result, such approaches may partially undermine the core principle of federated learning, namely the construction of a shared global representation that consistently integrates knowledge across all participating clients.

\textbf{Data Related Methods} To address label non-IID issues, another approach is to use dataset related information rather than modifying the model or the optimization process. Among them, FedLA leverages label-aware strategies to mitigate distribution skew; however, a key limitation of this method is its reliance on prior knowledge of the label distribution of each client, which must be available at initialization. This requirement raises significant privacy concerns and limits its applicability in realistic federated settings.
FedMcon proposes to alleviate data heterogeneity by exploiting contrastive objectives that rely on an additional dataset representative of the global data distribution. While effective in controlled environments, this assumption constitutes a concerning practical drawback, as such representative auxiliary datasets are never practially available in real-world federated learning applications. 


\section{Method}

\subsection{Privacy preserving aggregation}

Introdurre shamir secret sharing (almeno 2 immagini) per semplificare l'introduzione poi nel loop di aggregazione

\subsection{Problem Setting}

We consider a $k$-class classification problem over a dataset distributed across
multiple clients. Let $Y \in \mathbb{Z}_{\ge 0}^k$ denote the \emph{total label-count vector}
of the complete dataset, where $Y_\ell$ represents the total number of samples with
label $\ell \in \{1,\dots,k\}$, and let
\[
n \;=\; \sum_{\ell=1}^{k} Y_\ell
\]
The dataset is distributed across a set of clients
$\mathcal{C} = \{P_1,\dots,P_N\}$ participating in a federated learning task.
Each client $P_i$ holds a local dataset $\mathcal{S}_i$ of size $n_i = |\mathcal{S}_i|$.
For each client $P_i$, let
\[
Y_i = (c_{i,1},\dots,c_{i,k}) \in \mathbb{Z}_{\ge 0}^k
\]
denote the local label-count vector, where $c_{i,\ell}$ is the number of samples in
$\mathcal{S}_i$ with label $\ell$.

The goal is to train a global model $\theta \in \mathbb{R}^d$ via federated learning
in the presence of high statistical label heterogeneity across clients.
Unlike standard aggregation schemes, our method requires access to the
total sample distribution statistics to assign aggregation weights based on the client share of each label.

At communication round $r$, each participating client $P_i$ computes an updated local model $\theta_i^{r}$. The server aggregates the locally trained model as
\[
\theta^{r+1}
=
\theta^{r}
+
\sum_{i \in \mathcal{P}^{r}} w_i^{r} \, \theta_i^{r}
\]
where $\mathcal{P}^{r} \subseteq \mathcal{C}$ is the set of active clients at round $r$
and the weights $w_i^{r}$ are a function of sample distribution statistics as defined in X.X .

\subsection{Name of the Aggreation protocol}
\subsection{Client Routine}

At each communication round $t$, a subset of clients $\mathcal{P}^t \subseteq \mathcal{C}$
participates in the federated learning process.
Each participating client $P_i \in \mathcal{P}^t$ receives the current global model $\theta^{t-1}$ from the server and performs a local training phase using its local dataset
$\mathcal{S}_i$.

Starting from the received global model, client $P_i$ computes an updated local model $\theta_i^t$ by applying stochastic gradient descent on $\mathcal{S}_i$.
The local model update is then defined as the difference between the locally trained model and the global model,
\[
\Delta_i^t \leftarrow \theta_i^t - \theta^{t-1}.
\]

In addition to the model update, each client extracts statistics describing the distribution of labels in its local dataset.
In particular, client $P_i$ computes a label-count vector
\[
Y_i = (c_{i,1},\dots,c_{i,k}) \in \mathbb{Z}_{\ge 0}^k
\]
where each component of $Y_i$ represents the number of samples belonging to a given class.

To prevent disclosure of sensitive label information, the vector $Y_i$ is encoded using a $(t,n)$ Shamir secret sharing scheme over a finite field
$\mathbb{F}_p$, producing one secret share for each decryptor client in the set
$\mathcal{D} \subseteq \mathcal{C}$,
\[
\{\langle Y_i^t \rangle_j\}_{j \in \mathcal{D}}.
\]
Each share individually reveals no information about $Y_i^t$, while any subset of at least
$t$ shares allow its reconstruction.

Finally, client $P_i$ sends the model update $\Delta_i^t$ together with the set of secret shares $\{\langle Y_i^t \rangle_j\}_{j \in \mathcal{D}}$ to the server.
The server aggregates the received model updates according to the chosen federated learning strategy and keeps all the secret shares to the corresponding decryptor clients, enabling the secure computation of the aggregated label statistics as described in the Decryptor Routines.

\subsection{Decryptor Routines}
To enable the computation of the client's label statistics without revealing individual client
label information, we assume a subset $\mathcal{D} \subseteq \mathcal{C}$ of
decryptor clients with $|\mathcal{D}| = n$.
We work over a finite field $\mathbb{F}_p$ and adopt a $(t,n)$ Shamir secret sharing scheme,
where $t \le n$ is the reconstruction threshold.

During model initialization, each client $P_i$ encodes its label-count vector in $n$ secret shares
$\{\langle Y_i^t \rangle_j\}_{j \in \mathcal{D}} \in \mathbb{F}_p^{k}$ using a $(t,n)$ Shamir secret sharing scheme and sends the results to the server.
The server acts as a relay and at round $f$ distributes the received shares to the corresponding
decryptor clients $\{ P_j \}_{j \in \mathcal{D}}$. After receiving the share $\langle Y_i^t \rangle_j$, each decryptor aggregates the contribution received from the server according to
\[
\langle Y^{\mathrm{agg}} \rangle_j
\;\leftarrow\;
\langle Y^{\mathrm{agg}} \rangle_j
\;+\;
\sum_{i \in \mathcal{S}^t} \langle Y_i^t \rangle_j,
\]
where $\mathcal{S}^t$ denotes the set of different clients participating till round $t$.
This operation is performed locally, with low computational overhead, leveraging the linearity of Shamir's secret sharing method.
Once the aggregation is complete the aggregated shares $\{\langle Y^{\mathrm{agg}} \rangle_j\}_{j \in \mathcal{D}}$ are sent back to the server that is now capable of reconstructing the client label distribution.

\subsection{Server Algorithm}

The server never observes label statistics, neither in plaintext nor in shared form, and
only accesses the reconstructed aggregate when explicitly authorized by the decryptors.

\begin{algorithm}[h]
\caption{Federated Learning with Decryptor-Based Label Revelation}
\begin{algorithmic}[1]
\STATE Initialize global model $\theta^0$
\FOR{$t = 1$ to $T$}
    \STATE Sample client subset $S^t \subseteq \mathcal{C}$
    \FORALL{$i \in S^t$ \textbf{in parallel}}
        \STATE $\theta_i^t \leftarrow \texttt{ClientUpdate}(\theta^{t-1}, D_i)$
    \ENDFOR
    \STATE $\theta^t \leftarrow \texttt{AggregateModel}(\{\theta_i^t\}_{i \in S^t})$
    \IF{\texttt{TriggerCondition}(t)}
        \STATE Send reveal request to decryptors in $\mathcal{D}$
    \ENDIF
\ENDFOR
\end{algorithmic}
\end{algorithm}

\subsection{Hystogram Uncertainty Estimation}

\paragraph{Setup.}
Consider $K$ clients. Let $p_i \in [0,1]$ be the true proportion of label $\ell$
at client $i$, and define the global mean
\[
\mu = \frac{1}{K}\sum_{i=1}^K p_i .
\]
Each observed client samples with replacement $m$ examples,
\[
X_i \sim \mathrm{Binom}(m,p_i), 
\qquad 
\hat p_i = \frac{X_i}{m}.
\]
Let $K_{\text{obs}} = \alpha K$ be the number of observed clients and define
\[
\hat p = \frac{1}{K_{\text{obs}}}\sum_{i\in \mathrm{obs}} \hat p_i .
\]

\paragraph{Local variance (with replacement).}
\[
\mathbb{E}[\hat p_i \mid p_i] = p_i,
\qquad
\mathrm{Var}(\hat p_i \mid p_i) = \frac{1}{m} p_i(1-p_i).
\]

\paragraph{Variance decomposition.}
Using the law of total variance,
\[
\mathrm{Var}(\hat p-\mu)
=
\mathrm{Var}\!\left(\mathbb{E}[\hat p-\mu \mid \{p_i\},\mathrm{obs}]\right)
+
\mathbb{E}\!\left[\mathrm{Var}(\hat p-\mu \mid \{p_i\},\mathrm{obs})\right].
\]

\paragraph{Client sampling term.}
Define
\[
S^2 = \mathrm{Var}(p_i).
\]
If observed clients are sampled uniformly,
\[
\mathrm{Var}(\mu_{\mathrm{obs}}) 
= (1-\alpha)\frac{S^2}{K_{\mathrm{obs}}}.
\]

\paragraph{Local sampling term.}
\[
\mathbb{E}\!\left[\mathrm{Var}(\hat p \mid \{p_i\},\mathrm{obs})\right]
=
\frac{1}{K_{\mathrm{obs}}}\frac{1}{m}\,
\mathbb{E}[p_i(1-p_i)].
\]

\paragraph{Key identity.}
\[
\mathbb{E}[p_i(1-p_i)] 
= \mu(1-\mu) - S^2 .
\]

\paragraph{Total variance.}
\[
\mathrm{Var}(\hat p-\mu)
=
\frac{1}{K_{\mathrm{obs}}}
\left[
\left((1-\alpha)-\frac{1}{m}\right)S^2
+
\frac{1}{m}\mu(1-\mu)
\right].
\]

\paragraph{Normalized skew.}
Define
\[
\lambda = \frac{S^2}{\mu(1-\mu)} \in [0,1].
\]
Then
\[
\boxed{
\mathrm{Var}(\hat p-\mu)
=
\frac{\mu(1-\mu)}{K_{\mathrm{obs}}}
\left[
\frac{1}{m}
+
\lambda\!\left((1-\alpha)-\frac{1}{m}\right)
\right].
}
\]

\paragraph{Extremes.}
\[
\lambda=0 \Rightarrow 
\mathrm{Var}_{\mathrm{IID}}
=
\frac{\mu(1-\mu)}{K_{\mathrm{obs}}}\frac{1}{m},
\qquad
\lambda=1 \Rightarrow 
\mathrm{Var}_{\mathrm{max\ skew}}
=
\frac{\mu(1-\mu)}{K_{\mathrm{obs}}}(1-\alpha).
\]



\subsection{Security and Privacy Properties}

The proposed framework satisfies the following properties:

\begin{itemize}
    \item \textbf{Label privacy:}
    For any coalition of fewer than $t$ decryptors, the aggregated label statistics remain
    information-theoretically hidden.
    \item \textbf{Controlled disclosure:}
    Label distributions are revealed only upon explicit server request.
    \item \textbf{Low overhead:}
    Label aggregation relies exclusively on linear operations over small vectors.
    \item \textbf{Compatibility:}
    The approach is agnostic to the underlying federated optimization algorithm.
\end{itemize}


\section{Method}

\subsection{Overview}

We propose a federated learning framework that dynamically adapts aggregation strategies based on periodically estimated label distributions across clients. Our approach begins with standard FedAvg and transitions to a distribution-aware aggregation method after $r$ rounds. The labels histogram is estimated using a combination of Shamir secret sharing and distributed differential privacy, providing strong privacy guarantees tailored to medical imaging applications.

\subsection{Federated Learning Setup}

Let $\mathcal{C} = \{C_1, C_2, \ldots, C_N\}$ denote the set of $N$ participating clients, each possessing a private medical imaging dataset. A central server $S$ coordinates the training process over $T$ total rounds. We employ a hybrid aggregation strategy:
\begin{itemize}
    \item \textbf{Rounds $1$ to $r$}: Standard FedAvg aggregation
    \item \textbf{Rounds $kr+1$ onwards} (for $k = 1, 2, \ldots$): Distribution-aware aggregation using estimated label histograms.
\end{itemize}

The histogram estimation occurs each $kr$ rounds using federated analytics with cryptographic privacy guarantees.

\subsection{Privacy-Preserving Histogram Estimation}

\subsubsection{Protocol Architecture}

Our federated analytics protocol operates under an honest-but-curious security model for clients and decryptors, while the server may be malicious. The architecture consists of:
\begin{itemize}
    \item $N$ clients holding private label distributions
    \item A central server $S$ coordinating the computation
    \item A set of $M$ static decryptors $\mathcal{D} = \{D_1, D_2, \ldots, D_M\}$ with threshold $t = \lceil 2M/3 \rceil$ for consistency checking.
\end{itemize}

\subsubsection{Shamir Secret Sharing with Decryptors}

Each client $C_i$ possesses a local label histogram $Y_i = (c_{i,1},\dots,c_{i,k}) \in \mathbb{Z}_{\ge 0}^k$, where $K$ is the number of classes and each component of $Y_i$ represents the number of samples belonging to a given class. To securely aggregate these histograms without revealing individual contributions, we employ Shamir secret sharing~\cite{flamingo2023}:

\paragraph{Share Generation (One-time Setup):} At the beginning of training, each client $C_i$ generates a secret share per decryptor of its histogram:
\begin{enumerate}
    \item For each label count $c_{i,k}$ in  $Y_i$, client $C_i$ constructs a polynomial $f_{i,k}(x)$ of degree $t-1$ to encode the $c_{i,k}$ value.
    \item Create a share $s_{i,j}^{(k)} = f_{i,j}^{(k)}(x = {D}_k)$ for each decryptor $\mathcal{D}_k \in\{D_1, D_2, \ldots, D_M\}$
    \item Sends share $s_{i,j}^{(k)}$ to the server over a secure channel.
\end{enumerate}

The shares are transmitted only once at initialization, subsequent histogram estimations reuse these encrypted shares.

\paragraph{Server Buffering:} The server maintains a buffer $\mathcal{B}$ containing the identifiers of clients that have submitted shares. For histogram estimation at round $kr$ (where $k \geq 1$), the server determines the active client set $\mathcal{C}_{active} \subseteq \mathcal{C}$ participating in that round.

\subsubsection{Distributed Differential Privacy}

To ensure privacy against collusion between decryptors and the server, we apply distributed differential privacy using discrete Gaussian noise~\cite{kairouz2021distributed}. The noise addition is performed by the decryptors, preventing the server from learning exact histogram values.

\paragraph{Privacy Budget Allocation:} Let $\varepsilon_{total}$ be the total privacy budget across all $T$ rounds. For $K = \lfloor T/r \rfloor$ histogram estimations, each estimation receives budget $\varepsilon_r = \varepsilon_{total}/K$ under sequential composition.

\paragraph{Noise Generation:} At round $kr$, when histogram estimation is triggered:
\begin{enumerate}
    \item The server sends the secret shares encrypted from each client to the decryptor.
    \item Each decryptor $D_m$ independently computes the noise variance for round $kr$ as:
    \begin{equation}
        \sigma_m^2(kr) = \frac{L \cdot \Delta^2}{M \cdot |\mathcal{C}_{active}| \cdot \varepsilon_r^2}
    \end{equation}
    where $\Delta$ is the sensitivity (maximum change in histogram from one client 1/max number of samples per client), $L$ is the number of labels, and $M$ is the number of decryptors. This formulation ensures that the sum of noises from all $M$ decryptors provides $(\varepsilon_r, \delta)$-differential privacy~\cite{kairouz2021distributed,arxiv2601.09460}.
    
    \item Each decryptor $D_m$ samples discrete Gaussian noise $\mathbf{n}_m \sim \mathcal{N}_\mathbb{Z}^L(0, \sigma_m^2(kr))$ for each of the $L$ label counts
\end{enumerate}

Note that the noise variance is inversely proportional to $|\mathcal{C}_{active}|$, as specified in Equation 3 of~\cite{arxiv2601.09460}, ensuring proper calibration for the distributed setting.

\subsubsection{Secure Aggregation and Reconstruction}

At round $kr$, the histogram aggregation proceeds as follows:

\begin{enumerate}
    \item The server identifies active clients $\mathcal{C}_{active}$ and sends their identifiers to all decryptors
    
    \item Each decryptor $D_k$ performs:
    \begin{enumerate}
        \item Sums the shares: $\tilde{s}_j^{(k)} = \sum_{C_i \in \mathcal{C}_{active}} s_{i,j}^{(k)}$ for each label $j$
        \item Adds distributed noise: $\hat{s}_j^{(k)} = \tilde{s}_j^{(k)} + n_j^{(k)}$
        \item Sends noisy aggregated share $\hat{s}_j^{(k)}$ to the server
    \end{enumerate}
    
    \item \textbf{Consistency Check:} The server verifies that at least $t = \lceil 2M/3 \rceil$ decryptors provided valid shares. If fewer than $t$ decryptors respond or provide inconsistent shares, the estimation fails and FedAvg is used till another estimation is triggered.
    
    \item The server reconstructs the noisy histogram using Lagrange interpolation with any $t$ shares:
    \begin{equation}
        \hat{h}^{(j)} = \sum_{k=1}^{t} \hat{s}_j^{(m)} \prod_{\substack{\ell=1 \\ \ell \neq k}}^{t} \frac{-\ell}{k - \ell}
    \end{equation}
    
    \item The reconstructed histogram $\hat{\mathbf{h}} = [\hat{h}^{(1)}, \ldots, \hat{h}^{(L)}]$ is used to compute aggregation weights for the distribution-aware aggregation algorithm
\end{enumerate}

\subsection{Distribution-Aware Aggregation}

After obtaining the noisy histogram $\hat{\mathbf{h}}$ at round $kr$, we compute normalized class frequencies:
\begin{equation}
    \hat{p}_j = \frac{\max(0, \hat{h}^{(j)})}{\sum_{j'=1}^{L} \max(0, \hat{h}^{(j')})}
\end{equation}

These frequencies guide the aggregation strategy for rounds $(kr+1)$ through $((k+1)r)$. For instance, clients with underrepresented classes may receive higher weights, or class-specific regularization can be applied to address imbalance. The specific aggregation algorithm depends on the application and dataset characteristics.

\subsection{Security and Privacy Guarantees}

Our protocol provides the following guarantees:

\paragraph{Confidentiality:} Through Shamir secret sharing with threshold $t$, individual client histograms remain hidden as long as fewer than $t$ decryptors collude. With $t = \lceil 2M/3 \rceil$, the system tolerates up to $\lfloor M/3 \rfloor$ malicious decryptors~\cite{flamingo2023}.

\paragraph{Differential Privacy:} The distributed noise mechanism ensures $(\varepsilon_r, \delta)$-differential privacy for each histogram estimation, with total privacy budget $\varepsilon_{total}$ across all $K$ estimations through composition~\cite{kairouz2021distributed}. The discrete Gaussian noise prevents floating-point vulnerabilities common in continuous distributions.

\paragraph{Server Robustness:} By delegating noise addition to decryptors and requiring $t$-out-of-$M$ consistency, the protocol remains secure even if the server is malicious, as it cannot learn the true histogram without colluding with $t$ decryptors.

\subsection{Computational Complexity}

The protocol incurs minimal overhead:
\begin{itemize}
    \item \textbf{One-time share generation:} $O(NLM)$ for all clients to generate shares
    \item \textbf{Per estimation:} $O(|\mathcal{C}_{active}|LM)$ for share aggregation at decryptors and $O(tL)$ for reconstruction at the server
    \item \textbf{Communication:} Each client sends $O(LM)$ shares once; the server and decryptors exchange $O(LM)$ messages per estimation
\end{itemize}

Since histogram estimation occurs only every $r$ rounds, the amortized overhead per training round is negligible for moderate values of $r$.
\begin{table}
\caption{Test table.}\label{tab2}
%\scalebox{0.8}
{
\begin{tabular}{|l|l|l|}
\hline
Heading level &  Example & Font size and style\\
\hline
\small small &  {\Large\bfseries Lecture Notes} & 14 point, bold\\
\tiny tiny &  {\large\bfseries 1 Introduction} & 12 point, bold\\
\normalfont normal & {\bfseries 2.1 Printing Area} & 10 point, bold\\
\fontsize{8}{9} 8pt & {\bfseries Run-in Heading in Bold.} Text follows & 10 point, bold\\
\scriptsize script & {\itshape Lowest Level Heading.} Text follows & 10 point, italic\\
\hline
\end{tabular}}
\end{table}

\noindent Displayed equations are centered and set on a separate
line.
\begin{equation}
x + y = z
\end{equation}
Please try to avoid rasterized images for line-art diagrams and
schemas. Whenever possible, use vector graphics instead (see
Fig.~\ref{fig1}).

\begin{figure}
\includegraphics[width=\textwidth]{fig1.eps}
\caption{A figure caption is always placed below the illustration.
Please note that short captions are centered, while long ones are
justified by the macro package automatically.} \label{fig1}
\end{figure}

\begin{theorem}
This is a sample theorem. The run-in heading is set in bold, while
the following text appears in italics. Definitions, lemmas,
propositions, and corollaries are styled the same way.
\end{theorem}
%
% the environments 'definition', 'lemma', 'proposition', 'corollary',
% 'remark', and 'example' are defined in the LLNCS documentclass as well.
%
\begin{proof}
Proofs, examples, and remarks have the initial word in italics,
while the following text appears in normal font.
\end{proof}
For citations of references, we prefer the use of square brackets
and consecutive numbers. Citations using labels or the author/year
convention are also acceptable. The following bibliography provides
a sample reference list with entries for journal
articles~\cite{ref_article1}, an LNCS chapter~\cite{ref_lncs1}, a
book~\cite{ref_book1}, proceedings without editors~\cite{ref_proc1},
and a homepage~\cite{ref_url1}. Multiple citations are grouped
\cite{ref_article1,ref_lncs1,ref_book1},
\cite{ref_article1,ref_book1,ref_proc1,ref_url1}.

 %% removed for anonymized MICCAI submission.
    
    % The following acknowledgement and disclaimer sections can be removed for the double-blind review process.  If and when your paper is accepted, reinsert the acknowledgement and the disclaimer clause in your final camera-ready version.
    % IF you opted to include the acknowledgement and disclaimer sections, they will count towards the 8-page limit.

\begin{credits}
\subsubsection{\ackname} A bold run-in heading in small font size at the end of the paper is
used for general acknowledgments, for example: This study was funded
by X (grant number Y).

\subsubsection{\discintname}
It is now necessary to declare any competing interests or to specifically
state that the authors have no competing interests. Please place the
statement with a bold run-in heading in small font size beneath the
(optional) acknowledgments\footnote{If EquinOCS, our proceedings submission
system, is used, then the disclaimer can be provided directly in the system.},
for example: The authors have no competing interests to declare that are
relevant to the content of this article. Or: Author A has received research
grants from Company W. Author B has received a speaker honorarium from
Company X and owns stock in Company Y. Author C is a member of committee Z.
\end{credits}


%
% ---- Bibliography ----
%
% BibTeX users should specify bibliography style 'splncs04'.
% References will then be sorted and formatted in the correct style.
%
% \bibliographystyle{splncs04}
% \bibliography{mybibliography}
%
\begin{thebibliography}{8}
\bibitem{ref_article1}
Author, F.: Article title. Journal \textbf{2}(5), 99--110 (2016)

\bibitem{ref_lncs1}
Author, F., Author, S.: Title of a proceedings paper. In: Editor,
F., Editor, S. (eds.) CONFERENCE 2016, LNCS, vol. 9999, pp. 1--13.
Springer, Heidelberg (2016). \doi{10.10007/1234567890}

\bibitem{ref_book1}
Author, F., Author, S., Author, T.: Book title. 2nd edn. Publisher,
Location (1999)

\bibitem{ref_proc1}
Author, A.-B.: Contribution title. In: 9th International Proceedings
on Proceedings, pp. 1--2. Publisher, Location (2010)

\bibitem{ref_url1}
LNCS Homepage, \url{http://www.springer.com/lncs}, last accessed 2023/10/25
\end{thebibliography}
\end{document}
